\section{The Virgin and the Machine Gun}

\begin{quotex}
Un gentilhomme ne se moque pas de la religion des autres. \flright{\textsc{Cologero}}

Let satirists then laugh their fill at human affairs, let theologians rail, and let misanthropes praise to their utmost the life of untutored rusticity, let them heap contempt on men and praises on beasts; when all is said, they will find that men can provide for their wants much more easily by mutual help, and that only by uniting their forces can they escape from the dangers that on every side beset them: not to say how much more excellent and worthy of our knowledge it is, to study the actions of men than the actions of beasts. \flright{\textsc{Benedict Spinoza}}

The worst are full of passionate intensity. \flright{\textsc{W B Yeats}}

\end{quotex}
\paragraph{Evola the Dadaist}
There is a youtube video featuring Julius Evola explaining his early involvement with the Dada movement. The following is a translation from the French of the introduction to the video by its producers. 

\begin{quotationx}
These are some extracts of a televised interview with Julius Evola broadcast in 1971 by TFI, the French language Swiss television network. This interview brings to the views the period when Evola was a Dadaist painter.

In March 1971, I used to attend the School of Advanced Studies in Social Science in Paris to obtain a doctorate in political philosophy. But the cinema and television was already interesting me. One evening, I discussed with Jean-José Marchand who created ``The Archives of the 20th Century" for ORTF and this discussion led us to a fruitful collaboration. We were both animated with the desire to meet Julius Evola. We wanted to introduce him in a series of interviews bearing on three important points of Dadaism. I organized that interview which lasted a long time. At first, Evola was not entirely hostile but he remained skeptical. Then in impeccable French, he spoke to me for quite a while about his Dada experience and esoteric doctrines. The TV station retained just three minutes of that long dialog.

For posterity, I must point out that Evola refused to respond to two questions. The first:

``In the Livre du Gotha which belonged to my old comrade at the college in Geneva, Vittoria Emanuele de Savoie and his father Umberto, there is no ``Baron Evola" mentioned. Are you really a baron?"

The second:

``Why did you put a picture of Rudolf Steiner in your 1941 book, Sintesi di dottrina della razza, without mentioning his name? Yet you called him an example for the Nordic-Dinaric race, of the ascetic type, endowed with a power of spiritual penetration."

On that day, I understood that Steiner no longer interested him, indeed Evola did not like him. Evola impressed me greatly. 

\end{quotationx}
It's too bad that most the interview was discarded, particularly the much more important part about esoteric doctrines. Moreover, Evola's refusal to answer two questions is curious. As we've reported in these pages, a search through a history of Sicilian nobility also failed to mention a ``Baron Evola". Hence, we can dispense with the title even if it upsets those who believe such an authority is important. Besides, if he is a baron, where does that leave you in the restored world of tradition? One of the serfs, perhaps?

I've been investigating the reason for the Steiner photograph myself for several years, but have never received a satisfactory answer. Evola often quoted Steiner in his early philosophical works and he knew several Anthroposophists from the Ur/Krur groups. Yet it is curious that even as late as 1941 that he would refer to Steiner in such positive terms. It is certainly unclear how much of Steiner's occult Christianity that Evola was willing to consider.

\paragraph{The Virago, the M4, and the Rosary}

\begin{quotex}
From the female perspective, a woman identifies and rejects the coward because the coward will not protect her. The coward will abandon her and their children. \flright{\textsc{Ann Barnhardt}}

A man's sexual choice is the result and the sum of his fundamental convictions. Tell me what a man finds sexually attractive and I will tell you his entire philosophy on life. \flright{\textsc{Ayn Rand}}

\end{quotex}
The latest in the line of attractive viragos is Ann Barnhardt, in the space previously defined by Ann Coulter, or perhaps even as far back as St Joan of Arc. They seem to be taking what they conceive to be the man's role, probably for the reason stated above about cowards. Formerly the proprietor of a successful commodity brokerage, she gave it all up and converted to Traditional Catholicism. She then embraced Holy Virginity and Voluntary Poverty. I'm not sure about obedience. On the one hand she justifies usury, but most of her blog deals with the manifold injustices, not to mention financial risks, resulting from the usurious practices of certain financial institutions.

She also has high regard for the philosophy of Ayn Rand, provided it is baptized:

\begin{quotex}
Rand was right — reason is our only absolute, because Reason is God Himself.  If one re-reads Rand making this simple, conceptual substitution, it will literally knock you to the floor. 

\end{quotex}
I wonder who else has be knocked to the floor by that substitution? She even insists that Pope John Paul II's lectures on the \emph{Theology of the Body} were derived from Rand's views of sexuality.

She populates her blog with winsome love songs and ``man food" recipes, while boasting of her housecleaning skills. It would seem that she is playing Dagny Taggart trying to lure her John Galt. If you think you are the one, you'd better man up first.

\paragraph{Dante as Social Barometer}
\textbf{Dante}'s \emph{Divine Comedy} can be used as a barometer to determine the spiritual state of the modern world. To do that, you check out the inhabitants of the various levels of hell. A traditional society will recognize their sins for what they are. However, in the modern world, the sins are first tolerated, then defended, and finally glamorized.

Among the sinners in the seventh circle of hell, there are the \textbf{suicides}, \textbf{sodomites}, \textbf{usurers}, and \textbf{blasphemers}. I don't think much needs to be said each of them. Obviously, this list is outrageous to the modern mind, and that includes ``Fox news" type conservatives, not to mention most of the several versions of the ``other" rights. We only need to look at the praise give to celebrity suicides within the past year as well as changes to the marriage laws in several nations. As for usury, the USA and most of the world runs on debt. Usury is a disorder because it is sterile. Normally, an enterprise makes money through productivity and creativity, with money as the medium of exchange. However, for the financial institutions, money itself becomes the product through various forms of derivatives and so on.

Now the circle is complete with the worldwide praise for blasphemy and blasphemers, now defended in an unintellectual and an emotive way as a fundamental human right. Now murderers also inhabit the seventh circle because blasphemy was understood as violence against God, a crime once considered to be as serious as murder against other men. This opinion is absurd to people today, although it is just as much a part of their history.

Satirists are praised as somehow ``necessary" to free speech. That is true if correctly and intelligently done as, for example, we find in \textbf{Jonathan Swift} or \textbf{Boccaccio}, neither of whom is beyond employing a certain vulgarity or even scatological humor. Nevertheless, that is never an end in itself, but part of an intellectual whole.

Here we reach an abyss between Tradition and the modern world that accepts no compromise. If Dante is truly a high initiate, then it would be correct to claim that the modern West has now reached the seventh circle of hell. With this as a barometer, you can speculate about possible futures if you like by extrapolating the eight and ninth circles (in the comments).

Yet the modern West sees these inhabitants as praiseworthy, not sinners. Hence, the past needs to be annihilated and expunged, just as surely as the Taliban destroyed the Buddhist sculptures in Afghanistan a decade ago.

\paragraph{Intellectual Courage}
Several years ago, I was tutoring the 14 year old son of a friend in maths. The topic was the number line, so there were examples from temperature scales, and so on … anything with a plus and a minus. One particular problem was to state the number of years between 2 BC and AD 1. Now the obvious answer is ``3". However, the correct response is ``2" because there is no year ``0" in the scheme.

He seemed to understand and I challenged him to explain that when the question came up in class. The next time I saw him, I asked what happened; he told me he remained mute and accepted the incorrect answer from the teacher in silence. In other words, he distrusted his own knowledge and was cowed by the authority of the teacher and the textbook, as well as the conformity of his peers.

That displays a lack of intellectual courage. Similarly, those who know and understand traditional doctrines need to be certain of the correct answer and be willing to apply it when and where necessary.



\flrightit{Posted on 2015-01-11 by Cologero }

\begin{center}* * *\end{center}

\begin{footnotesize}\begin{sffamily}



\texttt{JA on 2015-01-13 at 20:35 said: }

For the record, Evola's family were barons under the Hohenstaufen kingdom of Sicily and were deprived of their rank after the French deposed the Staufens…..which is why they're not in the almanach de gotha…..


\end{sffamily}\end{footnotesize}
